\documentclass[11pt]{article}
\usepackage[T1]{fontenc}
\usepackage{lmodern,amsmath,amsthm,amssymb,mathtools,booktabs,microtype}
\usepackage[margin=1.03in]{geometry}
\usepackage[dvipsnames]{xcolor}
\usepackage[colorlinks=true,urlcolor=MidnightBlue,linkcolor=MidnightBlue,citecolor=MidnightBlue]{hyperref}
\newtheorem{theorem}{Theorem}[section]
\newtheorem{lemma}[theorem]{Lemma}
\newtheorem{corollary}[theorem]{Corollary}
\newtheorem{proposition}[theorem]{Proposition}
\theoremstyle{remark}\newtheorem{remark}[theorem]{Remark}
\DeclareMathOperator{\ord}{ord}
\DeclareMathOperator{\diag}{diag}
\DeclareMathOperator{\lcm}{lcm}
\newcommand{\C}{\mathbb C}\newcommand{\R}{\mathbb R}
\newcommand{\Z}{\mathbb Z}
\newcommand{\repo}{\href{https://github.com/VladimirReshetnikov/ProveIt/tree/00cadb961bbe77c74e5d082c4948f8884d58a970/Analysis/GeneralizedRiemannHypothesis}{ProveIt, pinned revision}}
\hypersetup{pdftitle={Recovering the characters and the graph from a coupled Dirichlet spectrum},pdfsubject={Inverse spectral theory for coupled Dirichlet L-functions}}
\usepackage{fancyhdr}
\pagestyle{fancy}\fancyhf{}\fancyhead[L]{\small\itshape Coupled Dirichlet spectra}\fancyhead[R]{\small Inverse problems}\fancyfoot[C]{\thepage}
\setlength{\headheight}{15pt}
\setlength{\emergencystretch}{2em}
\title{Recovering the characters and the graph\\from a coupled Dirichlet spectrum}
\author{Research note for Vladimir Reshetnikov}
\date{27 November 2026}
\begin{document}
\maketitle
\begin{abstract}
A Hermitian matrix couples a list of completed Dirichlet functions through a determinant of their logarithmic derivatives. We prove that its real zero multiset determines the unknown primitive character multiset and, for distinct characters, every principal minor of the matrix. The data may lose finitely many multiplicities, and the competitor may have arbitrary conductor and no more vertices than the source. The vertex bound is sharp. With finitely many insertions as well as deletions in the full spectrum, even the vertex count can be unknown. On the positive half-line alone, $o(T)$ edits through height $T$ preserve the labels. Even the coupling polynomial is identifiable after $o(T/(\log T)^{n+1})$ edits. For a fixed list we give a uniform finite observation window, a terminating procedure to certify a window, a linear-algebra decoder for false observations, and Lipschitz stability under perturbation of the surviving ordinates. The real-spectrum assertions use the Dirichlet GRH input of the ProveIt project. The arithmetic separation and polynomial independence arguments are unconditional.
\end{abstract}

\section{The inverse question}
A weighted graph usually has many isospectral competitors. In this problem each vertex carries an arithmetic label, and an observer sees the zeros of one entire function built from the labels and the graph Laplacian. The question is which features of an unknown graph and its labels the observer can distinguish. Ordinates are signed real numbers and multiplicities count.

The article in \repo\ combines 184 reports. Its section on zeros, Euler products and prime races constructs the coupled function, proves its real-root property under the stated input, and recovers all principal minors for a known list of distinct nonprincipal characters of a common parity. Its final inverse problem asks about mixed parity, unknown labels, effective observation windows, decoding in polynomial time, and perturbed ordinates. We address those questions in the observation models specified below. We use the upstream Dirichlet theorem at the level it is used in the article. This is a human-readable proof. The upstream Lean proof has not been rerun here.

\paragraph{Definition.}
For a primitive character $\chi$ of conductor $q>1$, put $a_\chi=(1-\chi(-1))/2\in\{0,1\}$ and
\[
 \Lambda_\chi(s)=(q/\pi)^{(s+a_\chi)/2}
       \Gamma((s+a_\chi)/2)L(s,\chi),\qquad
 F_\chi(z)=\eta_\chi^{-1}\Lambda_\chi(1/2+iz),
\]
where $\eta_\chi$ is a square root of its unit-modulus root number. For the primitive character of conductor one put $F_1(z)=\xi(1/2+iz)$, with $\xi(s)=s(s-1)\pi^{-s/2}\Gamma(s/2)\zeta(s)/2$. The functional equation makes these functions real entire of order at most one. Write $g_\chi=F'_\chi/F_\chi$ with derivative in $z$. For an ordered list $\boldsymbol\chi$ of length $n\ge1$, set $F_i=F_{\chi_i}$ and, for Hermitian $A$, define
\begin{align}
 P_A(z)&=\left(\prod_{i=1}^n F_i(z)\right)
 \det\bigl(A+\diag(g_{\chi_1}(z),\ldots,g_{\chi_n}(z))\bigr). \label{eq:model}
\end{align}
Expansion by principal minors removes the apparent poles. Choosing the opposite square root can only change an overall sign. Write $\Delta_I(f,h)=\sum_{x\in I}(\ord_x f-\ord_x h)_+$ for nonzero entire functions, and $\ord_x f=0$ when $f(x)\ne0$. A finite one-sided discrepancy puts no upper bound on the number of extra zeros of $h$.

\begin{theorem}[Recovery with unknown character labels]\label{thm:overview}
Use the Dirichlet GRH input. Let $A,B$ be Hermitian couplings on two lists of primitive Dirichlet characters, of lengths $n_A,n_B$. If $n_B\le n_A$ and $\Delta_\R(P_A,P_B)<\infty$, their character multisets are equal. For distinct labels, after identifying vertices by character, all corresponding principal minors are equal. In particular a nonnegatively weighted labelled graph is determined by these data, up to simultaneous relabelling. For repeated labels the exact invariant is the polynomial in \eqref{eq:group}. With finite discrepancy in both directions the size bound can be dropped.
\end{theorem}

\begin{theorem}[Recovery from positive ordinates alone]\label{thm:onehalf}
Under the same input, the positive zero multiset of an $n$-vertex coupling determines its primitive character multiset and its grouped determinant polynomial even after insertions and deletions whose total up to height $T$ is
\[
 o\!\left(\frac{T}{(\log T)^{n+1}}\right).
\]
The list length and the conductors can be unknown. In this statement each candidate must explain the data within the stated error rate for its own size. The labels alone are identifiable under the larger budget $o(T)$. Furthermore, if a competing list has size at most $n$, it suffices to bound the number of missing true positive ordinates by $o(T/(\log T)^{n+1})$; extra ordinates of the competitor can be unrestricted.
\end{theorem}
Theorem~\ref{thm:onehalf} follows from two different orders in the same spectral measurement. The leading average at a prime frequency gives the sum of character values at that prime. Once the labels agree, two different coupling polynomials have a first distinguishing frequency at a logarithmic order no greater than $n+1$. The proofs use prime-frequency traces and a positive trigonometric polynomial to count the two directions of discrepancy separately.

Here and throughout, graph weights refer to the off-diagonal weighted adjacency of an undirected loopless graph. Parallel edges, when allowed, contribute their total weight. Signed ordinates and exact complex characters matter; identifying the two ends of the real axis would change the observation model.

For instance the odd quadratic characters of conductors $3$ and $4$ and the even quadratic character of conductor $5$ can label a path with positive edge weights $a,b$. With these labels in path order its determinant polynomial is
\[
 X_1X_2X_3+aX_2X_3+(a+b)X_1X_3+bX_1X_2+ab(X_1+X_2+X_3).
\]
Theorem~\ref{thm:overview} identifies all three characters even against competing triples drawn from unbounded conductors. It identifies the middle character and both incident weights. The article already uses this particular list as a forward example. The inverse assertion for unknown inputs is the addition here.

\paragraph{Inputs and prior results.}
The project states its arithmetic input in \texttt{src:thm:dgrh}. Its forward result, \texttt{zer:cor:graphdirichlet}, bounds the difference between the output zero count and the sum of the input counts on any interval with regular endpoints by $n$. All zeros of the output are real under Dirichlet GRH. Indeed the genus-one factorization gives $\Im g_i(z)<0$ in the upper half-plane, since each $F_i$ has real zeros and at least one zero. The imaginary part of the diagonal matrix then rules out a null vector in \eqref{eq:model}, and conjugation handles the lower half-plane. Multiplicity and the count comparison are supplied by the article. The article's inverse argument is in \texttt{zer:thm:tailrigid} and \texttt{zer:thm:window}. Its prime-frequency reconstruction for one differentiated input is in \texttt{rec:thm:operatorinverse}. Background for the individual functions is in \cite{dlmf}.
Independence results for ordinary $L$-functions and derivatives, and uniqueness from arithmetic or zero data for $L$-functions themselves, have a substantial literature \cite{kmp,booker,nagoshi}. The determinant in \eqref{eq:model} generally has no given Euler product. Our use of an error locator is in the tradition of algebraic decoding \cite{pellikaan}. The contribution claimed here concerns these coupled spectra and the stated observation models. Historical priority has not been established exhaustively.

\section{Polynomial independence on a vertical line}
Let $\psi_1,\ldots,\psi_h$ be distinct primitive characters, including the conductor-one character if present. To avoid confusion with the digamma function, write $g_j=g_{\psi_j}$ and denote the digamma by $\Psi$. We first strengthen the analytic lemma underlying the article's fixed-list proof.
\begin{lemma}\label{lem:ap}
On $s=2+it$, $z=(s-1/2)/i$, a polynomial $P\in\C[X_1,\ldots,X_h]$ satisfying $P(g_1(z),\ldots,g_h(z))\to0$ as $t\to+\infty$ is zero. In particular the $g_j$ are algebraically independent over $\C(z)$.
\end{lemma}
\begin{proof}
The Euler product and the uniform digamma asymptotic give
\[
 g_j(z)=u(t)+c_j+D_j(t)+O((1+t)^{-1}),\quad
 u(t)=\tfrac i2\log((2+it)/2),\quad c_j=\tfrac i2\log(q_j/\pi),
\]
for $t\ge1$, with $D_j(t)=-i\sum_{m\ge2}\Lambda(m)\psi_j(m)m^{-2-it}$. Here $\Lambda(m)$ is the von Mangoldt function. The difference between $\Psi((s+a_j)/2)$ and $\log(s/2)$ is $O(1/t)$ for either parity. For $q_j=1$ the two extra logarithmic derivatives of $s(s-1)$ also have that size. Each $D_j$ is a uniform limit of trigonometric polynomials by absolute convergence.

For $d=\deg P$, replacing the $g_j$ by $u+c_j+D_j$ changes $P(g)$ by $o(1)$ since $|u|=O(\log(2+t))$. Expanding in $u$ gives $\sum_{k=0}^d u^k f_k(t)=o(1)$, where the $f_k$ are bounded uniform limits of trigonometric polynomials. Such a function with limit zero at $+\infty$ vanishes identically: simultaneous recurrence of the finitely many phases of a uniform approximant gives arbitrarily large translations arbitrarily close to periods. Equivalently one may take a sequence of translations approaching periods for successively finer approximants. Since $|u|\to\infty$, division by $u^d$ forces $f_d(t)\to0$, hence $f_d=0$. Descending induction gives $f_0=P(c_1+D_1,\ldots,c_h+D_h)=0$.

For $d>0$, let $\mathcal B$ polarize the top homogeneous term of $P$, with $\mathcal B(x,\ldots,x)=P_d(x)$. Extract the Fourier coefficient at $-\log m$ for $m=p_1\cdots p_d$ a product of distinct primes. Uniform absolute convergence permits termwise means. Only products from $P_d$ choosing one prime per factor can contribute, so the coefficient is
\[
 m^{-2}d!(-i)^d\prod_{k=1}^d\log p_k\,
 \mathcal B(v(p_1),\ldots,v(p_d)),\qquad v(p)=(\psi_j(p))_{j=1}^h.
\]
The distinct primitive characters induce distinct characters on the units modulo $Q=\lcm(q_1,\dots,q_h)$. Orthogonality shows that the corresponding vectors of character values span $\C^h$. Dirichlet's theorem on primes in progressions supplies distinct primes realizing any prescribed tuple of unit classes, including repeated classes. Multilinearity gives $\mathcal B=0$, a contradiction. A constant polynomial tending to zero is zero.

Finally, clear denominators of a putative rational-coefficient relation and write it as $\sum_{j=0}^b z^j P_j(g)=0$ with $P_b\ne0$. On our line $z=t-3i/2$ and every lower term after division by $z^b$ is $O(t^{-1}(1+\log t)^D)$ for some fixed $D$. Thus $P_b(g)\to0$, which is impossible.
\end{proof}

\section{Fixed labels, including repetitions}
If $m_j$ vertices carry label $\psi_j$, put $m=(m_1,\ldots,m_h)$ and define
\begin{equation}\label{eq:group}
 D_A(X)=\det\bigl(A+\diag(X_{j(i)})_{i=1}^{n}\bigr)
 =\sum_{0\le\alpha\le m} c_\alpha X^\alpha,\qquad c_m=1 .
\end{equation}
Here $j(i)$ identifies the character at vertex $i$ and inequalities on indices are componentwise. Each coefficient is the sum of the principal minors having that exponent pattern in the determinant expansion. Write
\[
 E_\alpha=\prod_{j=1}^h(F'_{\psi_j})^{\alpha_j}F_{\psi_j}^{m_j-\alpha_j},
 \qquad P_A=\sum_\alpha c_\alpha E_\alpha,\qquad
 V_\ell=\left\{\sum_{\alpha\ne m} q_\alpha(z)E_\alpha(z):\deg q_\alpha\le\ell\right\}.
\]
Lemma~\ref{lem:ap} makes this representation unique, also when polynomial multiples of $E_m$ are included. If all labels coincide, $D_A(X)=\det(A+XI)$, so even an exact full spectrum retains just the ordinary eigenvalue multiset of $A$.

\begin{theorem}[Polynomially weighted separation]\label{thm:window}
Under Dirichlet GRH, for every Hermitian $A$ and $0\ne H\in V_\ell$ one has $\Delta_\R(P_A,H)=\infty$. For a fixed list, $M\ge0$ and integers $K,\ell\ge0$, a finite $T_0$ works simultaneously for $\|A\|\le M$ and $0\ne H\in V_\ell$:
\[
 \Delta_{[-T_0,T_0]}(P_A,H)>K .
\]
In particular, for the same character list, $\Delta_\R(P_A,P_B)<\infty$ is equivalent to $D_A=D_B$. The forward implication also follows from discrepancy at most $K$ in the stated window for $\ell=0$, without a norm bound on $B$.
\end{theorem}
\begin{proof}
For $\Im z\le-3/2$ one has $\Re s\ge2$. Uniform Stirling asymptotics and bounded logarithmic derivative series there give $g_j(z)=\tfrac i2\log(s/2)+O(1)$ for $|s|\to\infty$. All terms in $D_A$ below $X^m$ have total degree at most $n-1$. Thus its leading value is $(\tfrac i2\log(s/2))^n(1+O(1/\log |s|))$. The quotient $H/P_A$ is $O((1+|z|)^\ell/\log |z|)$ there. Conjugation gives the same bound on $\Im z\ge3/2$, by conjugating the coefficients of $H$ when necessary.

Suppose there are finitely many poles. After multiplying by a polynomial $W$ cancelling their orders, $J=WH/P_A$ is entire of order at most one. Indeed the numerator and denominator are entire of order at most one, and their genus-at-most-one Hadamard products cancel, leaving a subproduct of the numerator times a linear exponential and a possible power of $z$. Equivalently the standard characteristic inequality for a quotient gives the same finite-order bound. Outside the strip, $J=O((1+|z|)^b)$ with $b=\deg W+\ell$. To propagate the estimate across the strip divide by $(z-2i)^b$ there and multiply by $\exp(-\delta\cosh(\kappa z))$ for $0<\kappa<\pi/3$, $\delta>0$. On long vertical sides this factor dominates finite-order growth, and its modulus is at most one on the horizontal boundaries. The maximum principle followed by $\delta\downarrow0$ proves the global polynomial bound. Cauchy estimates give that $J$ is a polynomial. Hence $H/P_A$ is rational. Dividing the identity $H=(H/P_A)P_A$ by $\prod_i F_i$ contradicts Lemma~\ref{lem:ap}: the coefficient of $X^m$ on the left is zero, and on the right it is the nonzero rational quotient. All zeros of $P_A$ are real, so the infinitely many poles prove the first claim.

If no finite window worked, take $A_j$ with $\|A_j\|\le M$ and $H_j\in V_\ell$ with unit coefficient vector, with discrepancy at most $K$ on $[-j,j]$. By finite-dimensional compactness pass to $A_j\to A$, $H_j\to H$ locally uniformly. Independence gives $H\ne0$ and $P_A\ne0$. By the first assertion, finitely many zeros of $P_A$ carry excess over $H$ exceeding $K$. Take disjoint small discs about them, enclosing no other zero of either limit, with boundaries free of zeros. Rouch\'e's theorem fixes the counts for $P_{A_j}$ and $H_j$ on these discs for all large $j$. Every zero of $P_{A_j}$ there is real, and the sum of positive differences of the disc counts is a lower bound for the discrepancy on $[-j,j]$. This contradicts the choice. Finally $P_B-P_A\in V_0$ vanishes to at least the common multiplicity at every zero. Applying separation to this difference, if nonzero, proves the last assertion.
\end{proof}
\section{Recovering the character labels}
\label{sec:blind}

Suppose the character list is itself unknown. A coupling of larger size can contain a smaller coupling as a direct summand, inheriting every one of its zeros. At equal or smaller size, even a one-sided comparison determines all the labels. We allow repetitions and include the primitive character of conductor one, with completion $\xi(s)=\tfrac12 s(s-1)\pi^{-s/2}\Gamma(s/2)\zeta(s)$.

\begin{theorem}[Blind recovery from a one-sided tail]\label{thm:blind}
Assume the Dirichlet critical-line input for the characters in two lists $\mathcal A,\mathcal B$ of sizes $n_A,n_B$. Let $A,B$ be Hermitian couplings on those lists. Suppose that all but finitely many zero multiplicities of $P_A$ occur among those of $P_B$. Then $n_B\ge n_A$. If $n_B=n_A$, the primitive character multisets coincide. After relabelling, the two determinant polynomials in the distinct label variables coincide, including all their grouped principal-minor coefficients. In particular, distinct labels determine the individual principal minors, and a nonnegatively weighted graph Laplacian is determined up to simultaneous relabelling when the labels are distinct.
\end{theorem}

\begin{proof}
The real-zero count of a Hermitian coupling on an interval with regular endpoints differs from the sum of the input counts by at most its size, by the inertia identity for the coupled determinant. The two-sided Riemann--von Mangoldt count for each fixed primitive completion on $[-T,T]$ is $(T/\pi)\log T+O_{\chi}(T)$. Passing to arbitrarily large regular endpoints gives $n_B\ge n_A$. Write $n=n_A=n_B$ for the rest of the proof.

Put $s=\tfrac12+iz$ and $C_A(s)=P_A((s-\tfrac12)/i)$, and define $C_B$ in the same way. For a list $\mathcal L$ let $Q_{\mathcal L}$ be the product of its conductors, $a_{\mathcal L}$ the sum of its parities (each $0$ or $1$), and $h_{\mathcal L}$ the number of copies of the conductor-one label. Define
\[
 c=\tfrac12\log(Q_{\mathcal B}/Q_{\mathcal A}),\qquad
 \beta=\tfrac12(a_{\mathcal B}-a_{\mathcal A})+2(h_{\mathcal B}-h_{\mathcal A}),\qquad
 G(s)=\frac{\prod_{\chi\in\mathcal B} L(s,\chi)}{\prod_{\chi\in\mathcal A} L(s,\chi)}.
\]
Every zero of $C_A$ is on the critical line. The one-sided hypothesis therefore supplies a monic polynomial $W$ cancelling all the poles of $C_B/C_A$ with their orders. Write $w=\deg W$, $e=w+\beta$ and $E=W C_B/C_A$. This entire function has order at most one. Indeed, both coupled functions are finite sums of products of entire functions of order at most one and their derivatives. In their genus-one Hadamard representations the denominator zeros cancel against those of $W C_B$, leaving a linear exponential and a canonical product over a submultiset of the zeros of $W C_B$. The counting bound $O_\epsilon(r^{1+\epsilon})$ on this submultiset gives order at most one.

Uniformly for $\Re s\ge 2$ and $|s|\to\infty$, there is a fixed $d\ne0$ for which
\begin{equation}\label{eq:blind-asymptotic}
 E(s)=d\,e^{cs}s^e G(s)\bigl(1+O(1/\log |s|)\bigr),
\end{equation}
using the principal logarithm in this half-plane. To verify the formula, divide each determinant by the product of its completions. The resulting determinant has diagonal logarithmic derivatives $\tfrac i2\log(s/2)+O(1)$ plus the fixed matrix. Here $L'/L$ is bounded by absolute convergence, the digamma estimate is uniform, and the extra logarithmic derivatives $1/s+1/(s-1)$ for $\xi$ are bounded. The two determinants both have leading term $(\tfrac i2\log(s/2))^n$ with relative error $O(1/\log|s|)$. In the ratio of the completion products, the equal numbers of gamma factors cancel their common part, and $\Gamma((s+1)/2)/\Gamma(s/2)=(s/2)^{1/2}(1+O(1/|s|))$. The additional $s(s-1)$ at each conductor-one occurrence contributes two powers. The conductor exponentials supply $e^{cs}$, the monic polynomial contributes $s^w(1+O(1/|s|))$, and all nonzero constant factors are absorbed into $d$.

Write $\kappa=|c|$. Since $G$ is bounded on $\Re s\ge 2$, the right half-plane satisfies $|E(s)|\ll (1+|s|)^e e^{c\Re s}$, enlarging the constant on a compact set. Phase normalization makes $P_A$ and $P_B$ real entire in $z$, hence $|C_j(1-\overline s)|=|C_j(s)|$. Reflection gives $|E(s)|\ll (1+|s|)^e e^{c(1-\Re s)}$ on $\Re s\le -1$, since the modulus of a monic polynomial of degree $w$ is comparable to $(1+|s|)^w$ off a compact set. In the intervening strip consider $E(s)/(s+3)^e$ with the analytic branch positive on the positive real axis. This quotient is bounded on both vertical boundaries. For $0<\alpha<\pi/3$, apply the maximum principle on expanding rectangles to the quotient multiplied by $\exp[-\delta\cos(\alpha(s-1/2))]$. The real part of the cosine is at least $\cos(3\alpha/2)\cosh(\alpha\Im s)$ throughout the strip, and the finite order bound controls the horizontal sides. Letting $\delta\downarrow0$ and combining the three regions yields
\begin{equation}\label{eq:blind-growth}
 |E(s)|\le C(1+|s|)^e\exp(\kappa|\Re s|)\qquad(s\in\C).
\end{equation}
Retaining the signed power $e$ in this estimate will remove the polynomial growth when the function is translated.

Choose $T_j\to+\infty$ so that $p^{-iT_j}\to1$ for every prime $p$ and $e^{icT_j}\to1$. For any finite list of these phases, simultaneous recurrence follows from compactness of the torus (or simultaneous Dirichlet approximation). Taking successively longer finite lists gives the sequence. Absolute convergence gives $G(s+iT_j)\to G(s)$ locally uniformly for $\Re s>1$, since the Dirichlet coefficients of the quotient are absolutely summable with weights $m^{-\sigma}$ for every $\sigma>1$. By \eqref{eq:blind-growth}, the entire functions $T_j^{-e}E(s+iT_j)$ are locally bounded. A subsequence converges locally uniformly to an entire function $H$ with
\[
 |H(s)|\le C e^{\kappa |\Re s|},\qquad
 H(s)=d_0 e^{cs}G(s)\quad(\Re s\ge2),\qquad d_0=d\exp(i\pi e/2)\ne0.
\]
Indeed, $(1+|s+iT_j|)/T_j\to1$ for each fixed $s$, and the second assertion follows from \eqref{eq:blind-asymptotic}.

Write $G(s)=\sum_{m\ge1} b_m m^{-s}$. For fixed $m$ put $\lambda=c-\log m$. On $\Re s=2$ the vertical mean of $H(s)e^{-\lambda s}$ equals $d_0 b_m$ by absolute convergence and orthogonality of distinct exponential frequencies. The same mean on $\Re s=-u$ has the same value for every $u>0$. To see this, integrate around the rectangle with those two vertical sides and horizontal sides at $\pm iY$, and divide by $2iY$. The horizontal integrals are $O_u(1)$ because of the global vertical-line bound for $H$. Consequently
\[
 |d_0 b_m|\le C\exp((\kappa+\lambda)u).
\]
Letting $u\to\infty$ proves that $G$ is a Dirichlet polynomial supported on $m\le e^{c+|c|}=\max(1,Q_{\mathcal B}/Q_{\mathcal A})$.

In particular $b_p=0$ for all sufficiently large primes. On the other hand, the prime coefficient of the quotient is
\[
 b_p=\sum_{\chi\in\mathcal B}\chi(p)-\sum_{\chi\in\mathcal A}\chi(p).
\]
For every unit residue class modulo the least common multiple of the conductors, Dirichlet's theorem supplies arbitrarily large primes in the class. Thus the two character sums agree on the full unit group. Distinct primitive characters induce distinct characters of this group, so orthogonality identifies all multiplicities. After the labels have been matched, the fixed-list one-sided rigidity theorem gives equality of the grouped determinant polynomials. This proves the assertion.
\end{proof}

\begin{corollary}[Finite corruption and unknown size]\label{cor:blind-symmetric}
A coupled zero divisor with finitely many arbitrary insertions and deletions determines its size, its multiset of primitive labels, and its grouped determinant polynomial. For one-sided containment, the size condition is sharp. In every larger size one can preserve the entire smaller divisor and introduce extra labels by taking a block-diagonal coupling with the original matrix as one block.
\end{corollary}
\begin{proof}
Apply the size inequality in both directions, then the equal-size conclusion. The block-diagonal determinant and the product of its vertex completions both factor, giving the stated example.
\end{proof}

\begin{corollary}[Finite windows for a bounded catalogue]\label{cor:blind-finite}
Fix a finite catalogue of primitive characters, a maximum list size $N$, a norm bound $M$, and an integer erasure budget $K\ge0$. There is $T<\infty$ such that the following implication holds for all lists drawn from the catalogue with $n_B\le n_A\le N$ and matrices of norm at most $M$. If at most $K$ zero multiplicities of $P_A$ in $[-T,T]$ are missing from $P_B$, then the lists coincide as multisets and the grouped determinant polynomials agree under relabelling.
\end{corollary}
\begin{proof}
First exclude unequal label multisets. If no window excluded them uniformly, choose counterexamples on $[-j,j]$. Finitely many choices of lists and compactness of the two matrix balls give a subsequence with fixed distinct multisets and convergent matrices. Both limiting determinants are nonzero and all their zeros are real. Around any finite set of distinct limit zeros of the first determinant choose disjoint disks whose closures contain no other zeros of either limit. Choose boundaries free of zeros of both functions. Rouch\'e's theorem preserves the two individual zero counts in each disk along the subsequence. The sum of the positive parts of their differences is at most $K$, since the approximating zeros are real, lie in the observed window in these disks, and the actual missing count dominates this sum. Exhausting the limit zeros gives one-sided tail discrepancy at most $K$. Theorem~\ref{thm:blind} now contradicts the choice of lists. Once the multisets agree, use the fixed-list finite-window theorem, taking the maximum of its thresholds over the finitely many lists. A finite-search interpretation of a sufficient window is given below.
\end{proof}


\paragraph{An unconditional complex-divisor version.}
For lists of the same specified size, Theorem~\ref{thm:blind} holds without Dirichlet GRH if the hypothesis counts missing zeros throughout $\C$. Indeed that hypothesis directly supplies the pole-clearing polynomial $W$. From that point on the proof uses the functional equation, bounds in the right half-plane, and the Euler coefficients of the primitive inputs. In the fixed-list proof the only use of the critical-line input in the infinite-divisor argument was to infer poles from missing real zeros. This locates precisely where the change from complex zero data to signed real ordinates uses GRH.

\paragraph{Reading a graph from the recovered coefficients.}
For distinct labels, the monomial supported on the complement of $J$ in \eqref{eq:group} has coefficient $\det A[J]$. In a loopless weighted Laplacian $L$ the one-vertex minors give the weighted degrees $L_{ii}$. For $i\ne j$ the edge weight, with zero for an absent edge, is
\[
 w_{ij}=\sqrt{L_{ii}L_{jj}-\det L[\{i,j\}]} .
\]
This proves the graph assertion in Theorem~\ref{thm:overview}. For a general Hermitian matrix, diagonal unitary conjugation and entrywise conjugation already preserve all principal minors. Our invariant is their full vector.

\section{Arithmetic traces from positive ordinates}
\label{sec:prime-trace}

A second observation model gives a quantitative test for the character labels. Let $\nu_A(\gamma)=\ord_\gamma P_A$ and $S_{\mathcal A}(m)=\sum_{\chi\in\mathcal A}\chi(m)$, with repetitions counted. The next formula transfers the classical prime-power trace of Landau \cite{gonek} to the zeros of the coupled determinant. Its leading term is independent of the Hermitian matrix.

\begin{theorem}\label{thm:landau}
Under the Dirichlet critical-line input, for each fixed integer $m\ge2$ and $T\ge3$,
\[
 \sum_{0<\gamma\le T}\nu_A(\gamma)m^{i\gamma}
 =-\frac{T\Lambda(m)}{2\pi\sqrt m}\,S_{\mathcal A}(m)
  +O_{A,\mathcal A,m}\!\left(\frac{T}{\log T}\right).
\]
\end{theorem}
\begin{proof}
Use $C_A(s)=P_A((s-1/2)/i)$, $s=2+it$, $b=3/2$ and $K_b(x)=b/(b^2+x^2)$. The genus-one Hadamard product and the critical-line location of all zeros imply
\[
 \Re\frac{C_A'(2+it)}{C_A(2+it)}=c_0+\sum_\gamma\nu_A(\gamma)K_b(t-\gamma)
\]
for a real constant $c_0$. The real parts of the regularizing terms $1/(1/2+i\gamma)$ are summable. The coupled count comparison and the unit-interval zero count for each primitive completion give $\sum_{u\le\gamma\le u+1}\nu_A(\gamma)\ll_{A,\mathcal A}\log(2+|u|)$. Endpoints can be included by slightly enlarging the interval.

Put $\omega=\log m$. The Fourier integral of $K_b$ at $\omega$ is $\pi e^{-b\omega}$, by residues. Integrating the preceding identity against $e^{i\omega t}$ for $0\le t\le T$ gives
\begin{equation}\label{eq:poisson-trace}
 \int_0^T e^{i\omega t}\Re\frac{C_A'(2+it)}{C_A(2+it)}\,dt
 =\pi e^{-b\omega}\sum_{0<\gamma\le T}\nu_A(\gamma)e^{i\omega\gamma}
 +O_{A,\mathcal A,m}(\log^2(2T)).
\end{equation}
Indeed, for $0<\gamma\le T$ the removed kernel tails have mass $O((1+\gamma)^{-1}+(1+T-\gamma)^{-1})$. For exterior ordinates in $[-2T,2T]$ the retained integral is $O((1+\operatorname{dist}(\gamma,[0,T]))^{-1})$. Summing over unit intervals costs $O(\log^2(2T))$. Farther away, the retained mass is $O(T/\gamma^2)$ with total $O(\log(2T))$. These bounds also justify integrating the sum. The constant integrates to $O_m(1)$.

On the same vertical line for large $t$, logarithmically differentiate the product of completions and the determinant in \eqref{eq:model}. The determinant term is $\operatorname{tr}(M^{-1}\partial_s M)$, with $M=A+\diag(g_i)$. We have $M=(\tfrac i2\log(s/2))I+O(1)$ and $\partial_s M=O(1)$ by the differentiated Euler series and the differentiated Stirling estimate. Consequently
\[
 \frac{C'_A}{C_A}(s)=\frac{n_A}{2}\log(s/2)+c_1
 -\sum_{k\ge2}\Lambda(k)S_{\mathcal A}(k) k^{-2-it}
 +O_{A,\mathcal A}(1/\log t).
\]
The polynomial factors in $\xi$ contribute $O(1/t)$. The integral of the real part of the logarithmic main term times $e^{i\omega t}$ is $O_m(\log(2T))$ by integration by parts. In the real part of the Euler sum only the negative frequency $-\log m$ resonates. All other summands have bounded integrals after twisting, uniformly in $k$ for fixed $m$, since the nonzero values of $|\log m\pm\log k|$ have a positive lower bound. Absolute convergence yields a total error $O_{\mathcal A,m}(1)$ from these terms. The integral is therefore $-T\Lambda(m)S_{\mathcal A}(m)/(2m^2)+O_{A,\mathcal A,m}(T/\log T)$. Substitution into \eqref{eq:poisson-trace} proves the formula.
\end{proof}

\begin{corollary}[Sublinear corruption of the positive spectrum]\label{cor:sublinear}
The positive ordinates alone determine the primitive label multiset even after $o(T)$ insertions and deletions through height $T$. In particular, for any two Hermitian couplings, writing $D(T)=\sum_{0<\gamma\le T}|\nu_A(\gamma)-\nu_B(\gamma)|$, one has
\[
 \liminf_{T\to\infty}\frac{D(T)}{T}\ \ge\
 \sup_{p\text{ prime}}\frac{\log p}{2\pi\sqrt p}\,
       |S_{\mathcal A}(p)-S_{\mathcal B}(p)|.
\]
The bound is positive whenever the label multisets differ, regardless of their sizes and conductors.
\end{corollary}
\begin{proof}
Each changed multiplicity alters a trace sum by a quantity of modulus at most one. Sublinear corruption preserves every normalized trace limit. Subtract the two trace formulas at a fixed prime, take absolute values and then take the supremum. If all prime contrasts vanish, use primes in each unit residue class of a common modulus and character orthogonality, as in Theorem~\ref{thm:blind}.
\end{proof}
The comparison here counts both excesses and deficits on the positive half-line. Its leading coefficients recover the labels. The coupling coefficients are recovered by the separate inverse theorems above.

\section{The matrix in lower-order positive-ordinate traces}
\label{sec:minor-trace}

The prime-power main term can agree for every frequency because the two character lists coincide. At a lower logarithmic order, a suitable frequency distinguishes every pair of distinct grouped determinant polynomials. This gives recovery from the positive half of the spectrum with a growing erasure budget.

\begin{theorem}\label{thm:minortrace}
Assume the critical-line input. Suppose $A,B$ have the same character multiset of size $n$, with $D_A\ne D_B$ after matching labels. There exist an integer $m\ge2$, an integer $1\le k\le n+1$, and $c_m\in\C\setminus\{0\}$ such that
\[
 \sum_{0<\gamma\le T}\bigl(\nu_A(\gamma)-\nu_B(\gamma)\bigr)m^{i\gamma}
 =c_m\frac{T}{(\log T)^k}
   +O\!\left(\frac{T}{(\log T)^{k+1}}\right).
\]
The implied constant may depend on both matrices, their labels, and $m$. Each of $\Delta_{(0,T]}(P_A,P_B)$ and $\Delta_{(0,T]}(P_B,P_A)$ is $\gg T/(\log T)^k$ for all sufficiently large $T$.
\end{theorem}

\begin{proof}
Let $w=(w_1,\ldots,w_h)$ have one indeterminate per distinct label, and let $\mathbf1$ be the all-ones vector. The two polynomials $D_A(U\mathbf1+w),D_B(U\mathbf1+w)$ are monic of degree $n$ in $U$. Expand formally at infinity, choosing the logarithm with constant term zero,
\[
 \log\frac{D_A(U\mathbf1+w)}{D_B(U\mathbf1+w)}
 =\sum_{r\ge1} h_r(w)U^{-r} .
\]
The $h_r$ are polynomials. The first nonzero $h_j$ has $j\le n$ since the two monic degree-$n$ polynomials differ (evaluate at $U=0$) and their difference has degree at most $n-1$. If $h_j$ is nonconstant, set $k=j$. Otherwise $h_j=c\ne0$, and translation of all the variables by $v$ replaces $U$ by $U+v$. Comparison at order $j+1$ gives $h_{j+1}(w+v\mathbf1)=h_{j+1}(w)-jc v$. In this case set $k=j+1$. Thus $k\le n+1$, $h_k$ is nonconstant, and all $h_r$ with $r<k$ are constant.

For $s=2+it$ put $u(s)=\tfrac i2\log(s/2)$ and let $w_j(s)=c_j-i\sum_{l\ge2}\Lambda(l)\psi_j(l)l^{-s}$ as in Lemma~\ref{lem:ap}. In particular $w,w'$ are bounded. The vector of actual logarithmic derivatives is $u\mathbf1+w+\epsilon$, with $\epsilon=O(t^{-1})$ and $\epsilon'=O(t^{-1})$, primes denoting $s$-derivatives. The first estimate holds uniformly in a fixed small disk about each $s=2+it$ by Stirling, which gives the second estimate by Cauchy's formula. The same estimates include the conductor-one correction. For bounded $w$ the formal logarithm above is an analytic power series in $U^{-1}$ on a disk of positive radius independent of $w$. Indeed the two monic polynomials divided by $U^n$ are uniformly close to $1$ there. This also holds on an enlarged bounded set of $w$, so the degree-$k$ truncation has remainder $O(|U|^{-k-1})$, with the same bound for its $w$-derivatives and $O(|U|^{-k-2})$ for its $U$-derivative. Substitute $U=u$, $w+\epsilon$ for $w$, and differentiate. Since all lower $h_r$ are constant, $u'=O(t^{-1})$, and the completion products cancel, the result is
\[
 \frac{C'_A}{C_A}(s)-\frac{C'_B}{C_B}(s)
 = u(s)^{-k}L_0(s)+O((\log t)^{-k-1}),\qquad
 L_0(s)=\frac{d}{ds}h_k(w(s)).
\]
No choice of logarithm of either coupled function away from this half-plane enters the formula.

The series $L_0(s)=\sum_{l\ge2} a_l l^{-s}$ converges absolutely at $\Re s=2$. Polynomial products of the $w_j$ converge absolutely at every real part greater than one, and inserting a factor $\log l$ in differentiating preserves convergence at $2$. At least one $a_m$ is nonzero. Otherwise $h_k(w(s))$ is constant, contrary to the squarefree-prime coefficient argument in Lemma~\ref{lem:ap}: subtract the putative constant and apply that argument directly to the nonconstant polynomial evaluated at $c_j-i\sum\Lambda(l)\psi_j(l)l^{-s}$. Fix such an $m$. For $t\ge3$ write $v(t)=(\log t)^{-k}$. Then $u^{-k}=(2/i)^k v(t)+O((\log t)^{-k-1})$. In the real part of $(2/i)^k L_0(2+it)$ the coefficient at frequency $-\log m$ is $(2/i)^k a_m m^{-2}/2$. The conjugate series has only positive frequencies. Each nonresonant integral $\int_3^T v(t)e^{i\lambda t}\,dt$ is $O_k(1/|\lambda|)$ by integration by parts; all nonzero frequency gaps for fixed $m$ are bounded below. The resonant integral is $T/(\log T)^k+O_k(T/(\log T)^{k+1})$ by integration by parts. Absolute summability and the preceding differential estimate now give
\[
 \int_0^T e^{it\log m}\Re\left(\frac{C'_A}{C_A}-\frac{C'_B}{C_B}\right)(2+it)\,dt
 =\frac{(2/i)^k a_m}{2m^2}\frac{T}{(\log T)^k}
   +O\!\left(\frac{T}{(\log T)^{k+1}}\right).
\]
Use \eqref{eq:poisson-trace} twice. Its $O(\log^2(2T))$ error is absorbed, and $c_m=(2/i)^k a_m/(2\pi\sqrt m)\ne0$. By the triangle inequality the sum of the two one-sided discrepancies is at least the modulus of the trace difference. Their difference is bounded: subtract the two coupled count comparisons with the same input counts, moving endpoints slightly if they are zeros so as to count exactly $(0,T]$. This proves the separate lower bounds.
\end{proof}

\begin{corollary}[Blind recovery from positive ordinates]\label{cor:blind-positive}
Suppose a positive-ordinate multiset is obtained from a coupling with $n$ primitive labels by insertions and deletions whose total through height $T$ is $o(T/(\log T)^{n+1})$. Among couplings satisfying this observation rule with their own vertex count, its character multiset and grouped determinant polynomial are unique. In particular, when all labels are distinct, the positive ordinates with these errors determine the unknown labelled graph Laplacian.
\end{corollary}
\begin{proof}
Two candidates first have symmetric discrepancy $o(T)$, so Corollary~\ref{cor:sublinear} identifies their labels and hence their common $n$. Their mutual discrepancy is then $o(T/(\log T)^{n+1})$. Theorem~\ref{thm:minortrace} forces equality of the grouped polynomials.
\end{proof}
For a known list the same reasoning gives a one-sided statement: $o(T/(\log T)^{n+1})$ missing multiplicities in $(0,T]$ force equality of the grouped polynomials even with unrestricted excess multiplicities. The exponent here comes from a fixed-frequency trace. The best possible discrepancy lower bound for the zero sets remains open.

\paragraph{Why order $n+1$ occurs.}
For $n\ge3$ give an $n$-cycle edge signs $+1$ and zero diagonal, and reverse just one edge sign for the second matrix. Write $A,B$ for the two signed adjacency matrices, with the same vertex labels. Every proper induced subgraph is a forest, on which a diagonal vertex sign change removes any prescribed edge signs. Its two principal minors therefore agree. In the full determinant, the only permutations affected are the two orientations of the cycle: traversing an edge as a transposition uses its sign twice. Hence
\[
 D_A(X)-D_B(X)=C,\qquad C=4(-1)^{n-1}.
\]
Writing $w_{j(i)}$ for the offset of vertex $i$, the coefficient of $U^{n-1}$ in $D_B(U\mathbf1+w)$ is $\sum_i w_{j(i)}$ since the diagonal of $B$ vanishes. Thus the first two nonzero orders of the formal logarithm are
\[
 CU^{-n}-C\left(\sum_i w_{j(i)}\right)U^{-(n+1)} .
\]
The constant coefficient at order $n$ contributes through $u'$ only, which is $O(1/t)$. The first nonconstant coefficient occurs at exactly $n+1$, even with distinct labels. This attains the last possible order of the frequency test in Theorem~\ref{thm:minortrace}. The number of individual zeros that differ could still be larger than the lower bound supplied by that test. The archive verifies the principal-minor identity in exact arithmetic for $3\le n\le8$.

For the triangle, attach the quadratic characters of fundamental discriminants $-3,-4,5$ in that order. At $p=2$ their values are $-1,0,-1$. Here $C=4$, $h_4=-4(w_1+w_2+w_3)$, and differentiation gives the Dirichlet coefficient $a_2=8i(\log 2)^2$. Thus the sign of a single cycle is visible in the explicit asymptotic
\[
 \sum_{0<\gamma\le T}(\nu_A(\gamma)-\nu_B(\gamma))\,2^{i\gamma}
 =\frac{64i(\log 2)^2}{\pi\sqrt2}\frac{T}{(\log T)^4}
   +O\!\left(\frac{T}{(\log T)^5}\right).
\]
The leading prime trace and the effects of every proper principal minor agree in this pair.

\section{Blind one-sided separation on the positive half-line}
\label{sec:riesz}

The article uses a positive prime-frequency weight for multiplicity differences of the input $L$-functions in \texttt{zer:thm:short}. Here the zeros belong to the coupled outputs. A single arithmetic trace supplies a lower bound for their total discrepancy. A nonnegative trigonometric polynomial combines several traces to force a linear discrepancy in each direction, even when one conductor product is much larger than the other.

\begin{theorem}\label{thm:directional}
Assume the critical-line input. If two couplings of the same size have different character multisets, there is $d>0$, depending on the pair, such that for all sufficiently large $T$,
\[
 \min\bigl\{\Delta_{(0,T]}(P_A,P_B),\Delta_{(0,T]}(P_B,P_A)\bigr\}\ge dT .
\]
\end{theorem}
\begin{proof}
For real $\omega$ put $J_T(\omega)=\sum_{0<\gamma\le T}(\nu_A(\gamma)-\nu_B(\gamma))e^{i\omega\gamma}$. For each fixed real $x>1$, the proof of Theorem~\ref{thm:landau} gives
\[
 J_T(\log x)=T b(x)+O_{A,B,\mathcal A,\mathcal B,x}(T/\log T),
 \quad
 b(x)=\begin{cases}
 -\dfrac{\log p}{2\pi\sqrt x}(S_{\mathcal A}(x)-S_{\mathcal B}(x)) &x=p^j,\quad j\ge1,\\[3pt]
 0 &\text{otherwise}.
 \end{cases}
\]
Indeed the Poisson argument applies to every $\log x>0$, and the distance from $\log x$ to the nearest different $\log l$, $l\ge2$ an integer, is positive for fixed $x$. The only possible resonant Euler term comes from $l=x$, with coefficient zero unless $x$ is a prime power. At the negative frequency $J_T(-\log x)=\overline{J_T(\log x)}$, so the main coefficient is $\overline{b(x)}$ and the same error estimate holds. The ordinary one-sided Riemann--von Mangoldt count for each primitive character is $T\log T/(2\pi)+O_\chi(T)$. Since the two lists have equal size, the coupled count comparison gives a constant $C_0$ with $|J_T(0)|\le C_0 T$ for $T\ge3$.

Choose a unit class $a$ modulo a common multiple of all conductors on which the two character sums differ. Such a class exists by character orthogonality. Along primes $p\equiv a$ the absolute value of that difference is a fixed positive number. The prime number theorem in this fixed progression gives $\sum_{p\equiv a}|b(p)|=\infty$; for example the primes in $[X,2X]$ already contribute $\gg\sqrt X$ for large $X$. Choose a finite set $\mathcal S$ of these primes with $\sum_{p\in\mathcal S}|b(p)|>C_0+2$, and choose $\theta_p$ with $e^{i\theta_p}b(p)=|b(p)|$. Consider the two Riesz products
\[
 W_\pm(t)=\prod_{p\in\mathcal S}\bigl(1\pm\cos(t\log p+\theta_p)\bigr),\qquad
 0\le W_\pm(t)\le 2^{|\mathcal S|}.
\]
Their constant terms are $1$ by unique factorization. Each term involving at least two distinct primes has frequency $\log r$, where $r$ is a positive rational with exponents $+1$ or $-1$ at two or more distinct primes. Neither $r$ nor $r^{-1}$ is a positive integral prime power. Consequently its trace has zero main term. The singleton terms contribute $\pm T|b(p)|$ for each prime, and the finite number of errors is $o(T)$. It follows that
\[
 \sum_{0<\gamma\le T}(\nu_A(\gamma)-\nu_B(\gamma))W_\pm(\gamma)
 =J_T(0)\ \pm T\sum_{p\in\mathcal S}|b(p)|+o(T).
\]
For large $T$ the plus expression is at least $T$ and the minus expression is at most $-T$. Nonnegativity and the upper bound on the weights give both claimed lower bounds, with $d=2^{-|\mathcal S|}$.
\end{proof}

\begin{corollary}[Positive ordinates, unrestricted surplus]\label{cor:optimal-heading}
Let the true list size be $n$. Suppose a competitor has at most $n$ vertices and the number of true positive-ordinate multiplicities missing from it through $T$ is $o(T/(\log T)^{n+1})$. Under the critical-line input, the character multisets and grouped determinant polynomials coincide. The competitor may have arbitrary conductor, matrix norm, and surplus multiplicities.
\end{corollary}
\begin{proof}
A smaller competitor is ruled out by the leading term of the ordinary positive-ordinate counts and the coupled count bound. At equal size Theorem~\ref{thm:directional} identifies the labels, then Theorem~\ref{thm:minortrace} identifies the grouped polynomial. A larger direct-sum competitor shows that the vertex restriction under one-sided observation is necessary.
\end{proof}

\section{Finite corrupted data and an exact decoder}
Assume the character multiset is known. Let $K$ bound deleted multiplicities and $\ell$ bound inserted multiplicities in the observation window. A spurious extra copy at a genuine zero counts as one insertion. For the window of Theorem~\ref{thm:window}, write the received multiset as $Y=Z+J$ where $Z$ consists of the true zeros after at most $K$ deletions and $|J|\le\ell$. The algorithm receives only $Y$. Denote the number of free determinant coefficients by $d=\prod_j(m_j+1)-1$.

\begin{theorem}[Exact decoding by linear equations]\label{thm:decode}
Seek a monic polynomial $W$ of degree $\ell$ and polynomials $Q_\alpha$ of degree at most $\ell$ such that
\[
 \left(W E_m+\sum_{\alpha\ne m} Q_\alpha E_\alpha\right)^{(r)}(y)=0
 \qquad (y\in\operatorname{supp}Y,\quad 0\le r<\operatorname{mult}_Y(y)).
\]
The system is consistent. In every solution $Q_\alpha=c_\alpha(A)W$, so the coefficient of $z^\ell$ in each $Q_\alpha$ gives the corresponding coefficient of $D_A$. No search over corrupted positions is required.
\end{theorem}
\begin{proof}
The choice $W=z^{\ell-|J|}\prod_{y\in J}(z-y)$, with multiplicities, and $Q_\alpha=c_\alpha W$ is feasible. In any solution subtract $WP_A$ from the fitted entire function. The difference is in $V_\ell$ and vanishes to the full retained orders along $Z$, since both terms do. Its missing multiplicity relative to $P_A$ is at most $K$ on the window. Hence Theorem~\ref{thm:window} forces the difference to vanish identically. Independence gives the polynomial identities coefficient by coefficient.
\end{proof}
There are $b=|Y|$ rows and $v=\ell+d(\ell+1)$ unknowns. Once exact jet values are available, matrix assembly by the product rule and Gaussian elimination take polynomially many arithmetic operations in $b,v$ (elimination alone takes $O((b+v)^3)$). The model allows exact real or complex arithmetic and zero tests. It makes no polynomial-bit-time assertion about computing or comparing transcendental ordinates. With distinct labels, $d=2^n-1$, the size of the requested vector of all minors. Extra capacity in the error budget may leave $W$ nonunique.

\paragraph{Checking a particular observation.}
The affine solution set is obtained by linear algebra. Take any solution, extract a candidate vector $c$ from the leading coefficients, and test the identities $Q_\alpha=c_\alpha W$ on the solution and the corresponding homogeneous identities on a basis of directions. If all pass, the candidate is certified under the given corruption model even in a shorter window, because the true parameters have at least one feasible lift. Enlarging the window with fresh complete observations subject to the same budgets eventually makes this test succeed.

\subsection*{A terminating search for a sufficient window}
There is also a uniform, computable height for a given fixed list, integer budgets and rational matrix bound $M$. The following construction gives an algorithm, potentially very slow. Its certificates use argument-principle counts on rational rectangles off the zeros. It does not need an oracle for exact real zero locations.

Parametrize Hermitian matrices by real coordinates and enclose their norm ball in a rational box of Hermitian matrices. Theorem~\ref{thm:window} applies on this enlarged compact set as well. For the coefficients of $H\in V_\ell$, normalize a nonzero coefficient of greatest modulus to one and take each remaining real and imaginary coordinate in $[-1,1]$. These finitely many compact rational boxes cover all scalar classes, and every point of every box gives a nonzero $H$. For each parameter pair $(A,H)$ there are finitely many disjoint rational rectangles whose two separate zero counts give
\[
 \sum_D \max\{N_D(P_A)-N_D(H),0\}>K .
\]
Indeed choose rectangles isolating the finite witnessing centers used in the compactness proof. Choose all boundaries free of zeros. The same counts persist in a parameter neighbourhood.

Here is explicit search logic.
\begin{enumerate}
\item Enumerate rational closed parameter boxes in the charts, finite disjoint rational rectangles in the complex plane, and rational boundary meshes and precision levels. Allow boxes to meet or extend across chart boundaries.
\item On rectangle boundaries use convergent complex interval enclosures of $P_A,H$ to certify, uniformly for a box, that each oriented mesh segment has image in an open disc excluding zero. At a rational representative parameter, compute the winding by a closed rational polygon: enclose its vertex values with sufficient precision that both endpoints of each interpolating segment belong to the same certified disc as the corresponding image arc. The polygon and image have the same winding by homotopy. Boundary nonvanishing throughout the box keeps it constant. Retain boxes for which the displayed inequality holds.
\item Dovetail these tests until the interiors of retained boxes cover each prescribed compact parameter chart. Coverage of rational boxes is decided by subdividing at their finitely many endpoints. Output any integer $T_0$ strictly greater than all absolute real parts of vertices of the rectangles used in this finite cover.
\end{enumerate}
The needed evaluations on compact sets are effective. For example one may use the standard rapidly convergent completed theta integrals split at $1$ via the functional equation, or the finite Hurwitz-zeta expression for $L$ with Euler--Maclaurin remainder bounds, together with analytic continuation across removable singularities by contour integration. Derivatives may likewise be enclosed by Cauchy's formula on a slightly larger disc. In enumerating rectangles one may avoid the discrete gamma poles on the boundaries. Equivalently interval evaluation on the edges may first subdivide into pole-free charts. The singular point of $\zeta$ can be avoided there as well. Absolute convergence with effective tail bounds gives arbitrarily fine enclosures on every such edge. Finite lists and arithmetic coefficients are effective inputs. At a witnessing parameter there is positive distance of the compact boundary images from zero. Sufficiently fine meshes and sufficiently small parameter boxes around that point therefore pass the tests. Compactness ensures termination. Reality of the zeros of $P_A$ converts the certified rectangle discrepancies into the required real-window bound. The same box search on pairs of bounded matrices for each pair of different multisets in a finite catalogue handles Corollary~\ref{cor:blind-finite}. A practical analytic bound on $T_0$ in terms of the conductors remains open.

\section{Stability under displacement and splitting of zeros}
\label{sec:stability}

Fix the labelled character data, a matrix norm bound $M$, an erasure budget $K$ and an insertion budget $\ell$. Write $m$ for the top multi-index, and let $d$ be the number of proper multi-indices $\alpha\ne m$ in the expansion
\[
 P_A=E_m+\sum_{\alpha\ne m} c_\alpha(A)E_\alpha.
\]
Use ordinary Euclidean norms on vectors of polynomial coefficients in increasing monomial order. The space $V_\ell$ consists of $\sum_{\alpha\ne m} h_\alpha E_\alpha$ with $\deg h_\alpha\le\ell$. Its coefficient norm is well-defined by rational linear independence. Let $T_0$ be a uniform window on which every nonzero member of $V_\ell$ misses more than $K$ zero multiplicities of every $P_A$ in the matrix ball, and put $T=T_0+1$. The margin permits roots at $\pm T$ to enter or leave the observation window as $A$ varies.

For a finite real multiset $D$ with $b$ elements counted with multiplicity, put $p_D(z)=\prod_{x\in D}(z-x)$. Write $I_D f$ for the unique polynomial of degree less than $b$ whose difference from the entire function $f$ is divisible by $p_D$ in the ring of entire functions. Thus a point of multiplicity $v$ prescribes derivatives of orders $0,\ldots,v-1$. Set $I_\varnothing f=0$. The coefficient norm of $I_Df$ treats nearby or coinciding nodes continuously.

\begin{lemma}\label{lem:stableinterpolation}
Fix a bound $B$ on $b$, and restrict all nodes to $[-U,U]$. If $D,D'$ have the same size and can be paired with distance at most $\delta$, then
\[
 \|I_Df-I_{D'}f\|\le C\delta\max_{|w|=U+2}|f(w)| .
\]
If $D\subseteq Y$ as multisets with $|Y|\le B$, then $\|I_D f\|\le C'\|I_Y f\|$. Both constants depend only on $B,U$.
\end{lemma}

\begin{proof}
On the positively oriented circle $|w|=U+2$, Cauchy's formula gives
\[
 I_D f(z)=\frac{1}{2\pi i}\int_{|w|=U+2}
 f(w)\frac{p_D(w)-p_D(z)}{(w-z)p_D(w)}\,dw .
\]
The right side is a polynomial in $z$ of degree less than $b$, and subtraction from $f(z)$ leaves a holomorphic multiple of $p_D(z)$ near the nodes. Each coefficient of its kernel has bounded first derivatives in the $b$ node variables on the compact box, since $|p_D(w)|\ge 2^b$. Taking maxima for $b\le B$ proves the first assertion, including collisions. For the second, $I_Df=I_D(I_Yf)$ follows from divisibility. The displayed integral is bounded uniformly on polynomials of degree less than $B$ with unit coefficient norm.
\end{proof}

\begin{lemma}\label{lem:stablefloor}
There is $a>0$ such that for all $\|A\|\le M$, every multiset $Z$ obtained from the zeros of $P_A$ in $[-T,T]$ by at most $K$ deletions, and every $H=\sum_{\alpha\ne m}h_\alpha E_\alpha\in V_\ell$, one has
\[
 \|I_ZH\|\ge a\,\|(h_\alpha)_{\alpha\ne m}\|.
\]
\end{lemma}

\begin{proof}
The number of zeros of $P_A$ in $[-T,T]$ is uniformly bounded. Indeed, local uniform convergence along any convergent subsequence of matrices gives a nonzero limit $P_{A_*}$. A circle enclosing the interval and avoiding its zeros then gives a common bound by Rouch\'e's theorem. If the claimed inequality failed, take a sequence $A_j\to A_*$, unit coefficient vectors converging to a unit vector defining $H_*$, and multisets $Z_j$ with $I_{Z_j}H_j\to0$. Pass to a subsequence of common cardinality with all ordered nodes convergent to a multiset $Z_*$. The integral formula gives $I_{Z_*}H_*=0$. Around each zero of $P_{A_*}$ in $[-T_0,T_0]$ take a disjoint small disk with its real diameter inside $(-T,T)$ and with boundary avoiding zeros of $P_{A_*}$ and points of $Z_*$. The exact zero counts of $P_{A_j}$ in the disks equal those of $P_{A_*}$ for large $j$, and all these zeros are real. Shrink the disks to exclude all points of $Z_*$ other than their respective centers. At most $K$ labels across these disks were deleted, so $Z_*$ contains all but at most $K$ multiplicities of the inner-window zeros of $P_{A_*}$. The divisibility $p_{Z_*}\mid H_*$ contradicts the choice of $T_0$.
\end{proof}

\begin{theorem}\label{thm:stabledecoder}
Suppose $\|A\|\le M$ and $0\le\delta\le1$. Erase at most $K$ labels of the zeros of $P_A$ in $[-T,T]$, displace each retained label by at most $\delta$, and insert at most $\ell$ arbitrary labels in $[-T-1,T+1]$. Let $Y$ be the resulting multiset, with no information about the matching. There are constants $R_0,C$ depending only on the fixed character list, $M,K,\ell,T$, such that the following estimator satisfies
\[
 \big\|(\widehat c_\alpha-c_\alpha(A))_{\alpha\ne m}\big\|\le C\delta .
\]
Choose any minimizer of
\[
 \left\| I_Y\!\left(W E_m+\sum_{\alpha\ne m}Q_\alpha E_\alpha\right)\right\|
\]
over monic $W$ of degree $\ell$ and $\deg Q_\alpha\le\ell$, with the joint coefficient norm bounded by $R_0$. Return $\widehat c_\alpha=[z^\ell]Q_\alpha$. The minimization is convex, with $\ell+d(\ell+1)$ free scalar coefficients. Repeated genuine zeros may split into separate reported locations, and distinct zeros may be reported at the same location.
\end{theorem}

\begin{proof}
Choose $R_0$ to contain all coefficient vectors
$(W_*,(c_\alpha(A)W_*)_\alpha)$ for $\|A\|\le M$ and all monic $W_*$ of degree $\ell$ whose roots lie in $[-T-1,T+1]$. Such a radius is finite by continuity and compactness. In particular, for the actual insertion multiset $J$, one may take $W_*=z^{\ell-|J|}p_J$. If $Z$ denotes the retained genuine labels before displacement, then $p_{Z\cup J}$ divides $W_*P_A$, where multiset union here adds multiplicities. The multiset $Z\cup J$ pairs with $Y$ at distance at most $\delta$. The first interpolation bound applied to the uniformly bounded family $W_*P_A$ shows that this feasible vector has residual at most $C_1\delta$. A minimizer exists by compactness and has at most this residual.

Write $G=W E_m+\sum Q_\alpha E_\alpha$ for its fitted function and $Z'$ for the displaced genuine submultiset of $Y$. The second interpolation bound gives $\|I_{Z'}G\|\le C_2\delta$. All fitted coefficient vectors lie in the same ball, so $|G|$ on the integration circle is uniformly bounded. The first bound then gives $\|I_ZG\|\le C_3\delta$. All cardinalities used here are bounded by the uniform count in the proof of the preceding lemma plus $\ell$. On $Z$ we have $I_Z(WP_A)=0$. Thus $H=G-WP_A=\sum (Q_\alpha-c_\alpha(A)W)E_\alpha$ satisfies the same residual bound there. Lemma~\ref{lem:stablefloor} yields
\[
 \big\|(Q_\alpha-c_\alpha(A)W)_{\alpha\ne m}\big\|\le (C_3/a)\delta .
\]
Since $W$ is monic, taking the coefficient of $z^\ell$ proves the result.
\end{proof}

The constants are existential. In particular, the theorem supplies no numerical condition number for a given list of characters. The contour representation makes the bound uniform over collisions; solving a system in unscaled point-evaluation rows can have very different numerical conditioning.


\section{Checks and remaining questions}
\paragraph{A mixed-parity example.}
The included computation uses the primitive quadratic character modulo $3$ (odd), the primitive quadratic character modulo $5$ (even), and the matrix
\[
 A=\begin{pmatrix}2&1\\1&3\end{pmatrix},\qquad
 D_A(X_1,X_2)=X_1X_2+3X_1+2X_2+5.
\]
At 45 decimal digits, the script evaluates the completions and their derivatives and refines nine sample zeros. It deletes the sample near $10.50530281214$ and inserts $4.25$. The linear system recovers the extra location $4.25$ and coefficients $3,2,5$. The normalized-row residual is about $2.4\times10^{-45}$. This experiment uses selected roots and numerical sign-change brackets. It does not certify completeness of an interval or a uniform window length. A separate exact rational polynomial experiment implements the jet conditions by polynomial divisibility. One copy is deleted, three are inserted, one genuine root is assigned excess multiplicity, and the resulting affine system has rank $15$. Its monic locator is $z(z+1)(z-87)$. It returns the same coefficient vector, with zero polynomial remainder. The polynomial example checks the decoder algebra in a finite model.

\paragraph{What remains open.}
Three questions are immediate. Can one bound the first sufficient uniform window explicitly by a useful function of the conductors, matrix norm and corruption budgets? Can one bound the stability constant at the same time? Can the graph be recovered in time polynomial in its vertex count when only the edge weights, rather than every principal minor, are wanted? For equal-size competitors the one-sided lower bound for different coupling polynomials is of order at least $T/(\log T)^{n+1}$. Do their actual positive zero multisets always differ on a larger scale, even for the signed cycles? The block-diagonal example describes a genuine obstruction when a larger competitor is admitted with unrestricted extra observations.

\begin{thebibliography}{9}
\bibitem{repo} ProveIt contributors, \emph{Beyond the Riemann hypothesis for Dirichlet $L$-functions: consequences, transfers and limits}, \repo, especially
\href{https://github.com/VladimirReshetnikov/ProveIt/blob/00cadb961bbe77c74e5d082c4948f8884d58a970/Analysis/GeneralizedRiemannHypothesis/Article/sections/zer.tex}{the section on zeros and prime races}. Upstream input: \href{https://github.com/lsernau/grh-lean/tree/72fead102f663d5c7a407964a13dc0cf9d977207}{\texttt{lsernau/grh-lean}, pinned revision}.
\bibitem{dlmf} NIST, \emph{Digital Library of Mathematical Functions}, \href{https://dlmf.nist.gov/25.15}{\S25.15} and \href{https://dlmf.nist.gov/5.11}{\S5.11}.
\bibitem{kmp} J. Kaczorowski, G. Molteni and A. Perelli, \emph{Linear independence of $L$-functions}, Forum Math. 18 (2006), 1--7. \href{https://sites.unimi.it/molteni/research/papers-pdf/14-molteni-Linear_independence_of_L-functions.pdf}{Author version}.
\bibitem{gonek} S. M. Gonek, \emph{A formula of Landau and mean values of $\zeta(s)$}, in \emph{Topics in Analytic Number Theory}, 92--97, 1985. \href{https://www.sas.rochester.edu/mth/people/faculty/gonek-steve/assets/pdf/8-landau-form.pdf}{Author version}.
\bibitem{booker} A. R. Booker, \emph{$L$-functions as distributions}, Math. Ann. 363 (2015), 423--454. \href{https://doi.org/10.1007/s00208-015-1178-z}{DOI}.
\bibitem{nagoshi} H. Nagoshi, \emph{Independence of $L$-functions}, \href{https://www.math.kobe-u.ac.jp/prob/nagoshi_1.pdf}{research abstract}. See also B. Bagchi, \emph{A joint universality theorem for Dirichlet $L$-functions}, Math. Z. 181 (1982), 319--334, \href{https://doi.org/10.1007/BF01161980}{DOI}.
\bibitem{pellikaan} R. Pellikaan, \emph{On decoding by error location and dependent sets of error positions}, Discrete Math. 106/107 (1992), 369--381. \href{https://doi.org/10.1016/0012-365X(92)90567-Y}{DOI}.
\end{thebibliography}
\end{document}
